\section{Methods} \label{methods}
{\color{red}

Below, we first outline the \texttt{gamdid} approach (model equations,
test statistics and the simultaneous confidence bands are discussed in
more detail in the Supplementary materials).
We subsequently describe the covariate adjustment procedure, the datasets,
the simulation framework and the implementation.

\subsection{Overview of the gamdid approach}

The analysis is performed for each feature (protein or peptide) separately and
consists of five steps (see Figure \ref{fig:workflow}).
\\
\\
\textbf{Step 1: Estimating a smooth distribution per group}
For a given feature, the abundance values of all cells are binned into a histogram,
using the same bins for every group (see Figure \ref{fig:workflow}, panel A).
Following Lindsey \cite{lindsay_method} and Efron \cite{efron}, the bin counts
can be treated as Poisson counts.
\\
We then model the bin counts as a function of the bin midpoints using a
Poisson GAM, which allows us to flexibly model the
shape of the distribution with a smooth curve. Each group receives its own
smooth curve (see Figure \ref{fig:workflow}, panel B).
Equation \ref{eq:GAM_simple} schematically represents the smooth curve for group $c$
and bin $k$ with midpoint $mid_k$.
\begin{equation}
    \log(\text{expected count}_{ck}) = \underbrace{f_c(mid_k)}_{\text{density estimation for group } c} + \text{normalization offset}_c
    \label{eq:GAM_simple}
\end{equation}
where $f_c$ is a flexible smooth function that captures the shape of the distribution of group $c$.
The counts are normalized using an offset, which allows us to interpret the curves
as densities.
\\
\\
\textbf{Step 2: Testing whether the distributions differ}
If two groups have the same distribution, their smooth curves are identical,
i.e. $H_0: f_{c_1} = f_{c_2}$. We test this by assessing if
the parameters corresponding to a smooth curve $\beta_{f_c}$
differ between the groups, i.e. $H_0: \beta_{f_{c_1}} = \beta_{f_{c_2}}$.
We do this by employing a Wald-type test that builds on the work of Wood \cite{Wald}.
The resulting p-value, referred to as the GAM p-value, captures any difference in
distribution between the groups: mean shifts, differences in variance, modality,
skewness...
\\
\\
\textbf{Step 3: Combining with a mean-shift test}
Because the GAM test spreads its power over all possible types of distributional
differences,
it is less powerful than a dedicated DA method for detecting pure mean shifts.
We therefore also compute a p-value with \texttt{msqrob2}, a robust linear-model
approach designed for detecting mean shifts in proteomics data \cite{msqrob}.
The GAM and
\texttt{msqrob2} p-values are combined into a single 'feature-level aggregated
p-value' using the harmonic mean p-value (HMP) method \cite{harmonicp}.
As a result, \texttt{gamdid} largely retains the
sensitivity of conventional DA analysis for mean shifts, while additionally
detecting changes in distributional shape that DA methods miss.
\\
\\
\textbf{Step 4: Multiple testing and pairwise comparisons}
The feature-level aggregated p-values are corrected for multiple testing across
features using the Benjamini-Hochberg FDR procedure
\cite{BH_procedure}. When more than two groups are compared, significant features
are followed up with pairwise comparisons between all pairs of groups, using the
same combination of GAM and \texttt{msqrob2} tests. The feature-level and pairwise
hypotheses are jointly controlled for multiple testing by the stagewise testing
procedure \cite{stagewise_testing, stager}.
\\
\\
\textbf{Step 5: Detecting differential regions}
Finally, \texttt{gamdid} can indicate where the
distributions differ. For this, a log-difference smoother is computed. This
smoother displays the difference between the estimated log-densities of two
groups, together with a simultaneous confidence band. Regions where this
band excludes zero indicate a statistically significant local distributional
difference; these are marked by gray shading (Figure \ref{fig:workflow}, panel C).
}

\subsection{Covariate adjustment}
In addition to the covariate of interest, \texttt{gamdid} allows users to account
for confounding variables through a nuisance formula, enabling applications such
as batch correction in multi-batch experimental designs. Both the covariate of
interest and the nuisance formula are combined into one model in the \texttt{msqrob2}
analysis, and an appropriate contrast is constructed to test for distributional
differences in the covariate of interest.\\
To remove the influence of nuisance covariates from the GAM-based inference,
partial residuals with respect to the covariate of interest are computed from
the \texttt{msqrob2} fit and subsequently used as input to the GAM analysis.
This two-step approach is required due to the binning strategy: as a single bin
may contain observations that correspond to different groups of the nuisance
covariate, direct covariate adjustment within the GAM framework is not feasible.\\
Note, however, that the uncertainty associated with the regression of the nuisance
covariates in the first step is not accounted for, so the resulting p-values should be
interpreted with caution.

\subsection{Data}
\begin{itemize}
    \item \texttt{leduc2022\_pSCoPE}: These data were acquired using pSCoPE technology: a TMT-18plex protocol with prioritized data-dependent acquisition (DDA) carried out by a Thermo Scientific Q-Exactive mass spectrometer \cite{leduc}. The MS data were identified and quantified using MaxQuant version 1.6.17 \cite{maxquant}. We performed quality control by filtering low-quality features and filtering cells with a median coefficient of variation greater than 0.6, calculated using the \textit{medianCVperCell} function of the \texttt{scp} package \cite{scp_paper}. The PSM intensities were log2-transformed for variance-stabilization, median centered and summarized to the protein level using median summarization per MS run. Subsequently, samples were normalized through median centering and proteins with more than 95\% missingness were filtered out, resulting in a dataset with 776 melanoma cells and 755 monocytes profiled across 2120 proteins.
    \item \texttt{petrosius2023\_mES}: These data originate from a label-free SCP experiment conducted on mouse embryonic stem cells cultured under two distinct conditions: a serum-free 2i condition (m2i) containing cytokine LIF and inhibitors targeting the MEK and GSK3 pathways, and a serum condition (m15) only containing cytokine LIF \cite{petro}. Data acquisition was performed using an Orbitrap Eclipse Tribrid mass spectrometer operating in data-independent acquisition (DIA) mode. The resulting mass spectrometry data were subsequently analyzed by the authors using Spectronaut v17 \cite{spectronaut}.
    We performed quality control by removing cells with fewer than 500 detected peptides or a median intensity lower than 11. Features with more than 95\% missing values were also removed. The detected intensities were log2-transformed for variance-stabilization, and the intensities of each cell were centered with the median log2-intensity of the corresponding cell. Subsequently, peptides were summarized to the protein level through median polish summarization. The resulting dataset contains 295 m15 and 269 m2i cells with 2492 proteins.
\end{itemize}
Preprocessed protein-level summarized assays from both datasets served as the foundation for our simulations. From the Petrosius dataset, we selected the 269 m2i-condition cells, while from Leduc, we selected the 755 monocytes (in both cases following removal of non-informative cells during preprocessing). Within each selected subset, no systematic biological differences are expected between cells. Consequently, random assignment of cells to two mock groups (A and B) yields a dataset suitable as a mock basis for a controlled two-group comparative analysis.

\subsection{Simulation framework}
The benchmark simulations are based on protein-level mock data from both the Leduc and Petrosius datasets, as described in the Results section. For each dataset, mock groups A and B were formed by random assignment of cells, and distributional signal was introduced in group B for 20\% of randomly selected features. Five types of distributional signal were considered, following the framework of Korthauer et al. (2016) \cite{scdd}:
\begin{itemize}
    \item \textbf{DE:} A shift equal to 40\% of the feature's standard deviation is added to all observations in group B.
    \item \textbf{DV:} Independent Gaussian noise with mean zero and standard deviation equal to 120\% of the feature's standard deviation is added to all observations in group B, thereby increasing the spread of the distribution without altering its mean.
    \item \textbf{DB:} For 50\% of the observations in group B, a value of 90\% of the feature's standard deviation is added; for the remaining 50\%, the same value is subtracted. This induces a bimodal distribution in group B.
    \item \textbf{DM:} A value of 90\% of the feature's standard deviation is added to 50\% of the observations in group B, creating an asymmetric shift concentrated in one tail.
    \item \textbf{DP:} In both groups, observations are split into two subpopulations that are shifted upwards or downwards by 100\% of the feature's standard deviation. In group A, 40\% of the observations are shifted upwards and the remaining 60\% downwards; in group B, 70\% are shifted upwards and 30\% downwards. Both groups are therefore bimodal with identical component locations but different mixing proportions, which also results in a difference in group means of 0.6 standard deviations. In the three-group simulations, group C is treated identically to group A.
\end{itemize}


\subsection{Implementation}
The methods for assessing differential distributions are implemented in the \texttt{gamdid} R package that is publicly available through GitHub (github.com/statOmics/gamdid) and will be submitted to Bioconductor.\\
\texttt{gamdid} builds upon the CRAN packages \texttt{mgcv} \cite{mgcv} and \texttt{harmonicmeanp} \cite{hmp_package} for fitting generalized additive models and p-value aggregation; and on the R/Bioconductor packages \texttt{stageR} \cite{stager} for stagewise testing and \texttt{msqrob2} for detecting shifts in mean abundance \cite{msqrob}. Specifically, \texttt{mgcv} version 1.9.4, \texttt{harmonicmeanp} version 3.0.1, \texttt{stageR} version 1.32.0 and \texttt{msqrob2} version 1.18.0 were used in R version 4.5.3.
